% Options for packages loaded elsewhere
\PassOptionsToPackage{unicode}{hyperref}
\PassOptionsToPackage{hyphens}{url}
\documentclass[
  a4paper,
]{article}
\usepackage{xcolor}
\usepackage[margin=25mm]{geometry}
\usepackage{amsmath,amssymb}
\setcounter{secnumdepth}{-\maxdimen} % remove section numbering
\usepackage{iftex}
\ifPDFTeX
  \usepackage[T1]{fontenc}
  \usepackage[utf8]{inputenc}
  \usepackage{textcomp} % provide euro and other symbols
\else % if luatex or xetex
  \usepackage{unicode-math} % this also loads fontspec
  \defaultfontfeatures{Scale=MatchLowercase}
  \defaultfontfeatures[\rmfamily]{Ligatures=TeX,Scale=1}
\fi
\usepackage{lmodern}
\ifPDFTeX\else
  % xetex/luatex font selection
\fi
% Use upquote if available, for straight quotes in verbatim environments
\IfFileExists{upquote.sty}{\usepackage{upquote}}{}
\IfFileExists{microtype.sty}{% use microtype if available
  \usepackage[]{microtype}
  \UseMicrotypeSet[protrusion]{basicmath} % disable protrusion for tt fonts
}{}
\makeatletter
\@ifundefined{KOMAClassName}{% if non-KOMA class
  \IfFileExists{parskip.sty}{%
    \usepackage{parskip}
  }{% else
    \setlength{\parindent}{0pt}
    \setlength{\parskip}{6pt plus 2pt minus 1pt}}
}{% if KOMA class
  \KOMAoptions{parskip=half}}
\makeatother
\usepackage{longtable,booktabs,array}
\newcounter{none} % for unnumbered tables
\usepackage{calc} % for calculating minipage widths
% Correct order of tables after \paragraph or \subparagraph
\usepackage{etoolbox}
\makeatletter
\patchcmd\longtable{\par}{\if@noskipsec\mbox{}\fi\par}{}{}
\makeatother
% Allow footnotes in longtable head/foot
\IfFileExists{footnotehyper.sty}{\usepackage{footnotehyper}}{\usepackage{footnote}}
\makesavenoteenv{longtable}
\usepackage{graphicx}
\makeatletter
\newsavebox\pandoc@box
\newcommand*\pandocbounded[1]{% scales image to fit in text height/width
  \sbox\pandoc@box{#1}%
  \Gscale@div\@tempa{\textheight}{\dimexpr\ht\pandoc@box+\dp\pandoc@box\relax}%
  \Gscale@div\@tempb{\linewidth}{\wd\pandoc@box}%
  \ifdim\@tempb\p@<\@tempa\p@\let\@tempa\@tempb\fi% select the smaller of both
  \ifdim\@tempa\p@<\p@\scalebox{\@tempa}{\usebox\pandoc@box}%
  \else\usebox{\pandoc@box}%
  \fi%
}
% Set default figure placement to htbp
\def\fps@figure{htbp}
\makeatother
\setlength{\emergencystretch}{3em} % prevent overfull lines
\providecommand{\tightlist}{%
  \setlength{\itemsep}{0pt}\setlength{\parskip}{0pt}}
\usepackage{bookmark}
\IfFileExists{xurl.sty}{\usepackage{xurl}}{} % add URL line breaks if available
\urlstyle{same}
\hypersetup{
  pdftitle={A Finite Complete-Relation Closed System with Time Evolution as Vertices: Exact Minimal Closure at N=1{,} K=4 under U\^{}K=I and Complete-Network Analysis of Even and Odd Harmonics},
  pdfauthor={Noriaki Kihara (WF System Co., Ltd.), ORCID 0009-0004-6753-4020},
  hidelinks,
  pdfcreator={LaTeX via pandoc}}

\title{A Finite Complete-Relation Closed System with Time Evolution as
Vertices: Exact Minimal Closure at \(N=1, K=4\) under \(U^K=I\) and
Complete-Network Analysis of Even and Odd Harmonics}
\author{Noriaki Kihara (WF System Co., Ltd.), ORCID 0009-0004-6753-4020}
\date{September 15, 2026. v1.2. Version DOI 10.5281/zenodo.22763330}

\begin{document}
\maketitle

\subsection{Abstract}\label{abstract}

The aim of this study is to analyze a minimal exactly closed system in
which time evolution is not given as an external distinguished
parameter: states at different phase instants are treated as ordinary
state vertices, without distinction, and all their relations are closed
into a complete network. The starting point is the numerical experiment
on a self-consistent closed system of relational waves in the preceding
study {[}K1{]}. There, \(N\) vertices and \(M=N(N-1)/2\) relational
waves were placed at a nominal instant, and the interaction was iterated
with respect to a discrete update variable \(\tau\) with
\(\Delta\tau=2\pi/N\). For highly symmetric initial states, the
relational waves were observed to concentrate in narrow bands near the
90-degree series in the complex plane --- a cross-shaped configuration
--- and then, after a rapid growth of the transverse component, to
transition to a quasi-stable ring of nearly equal amplitudes with widely
dispersed phases. However, this description singles out \(\tau\) as a
special update direction distinguished from the other state variables,
and presupposes ``\(N\) vertices at a given instant''; both points leave
a possible violation of the anonymity principle.

In this paper, when each nominal instant carries \(N\) states and the
phase closure carries \(K\) phase states, we treat all states without
distinction as \(P=KN\) vertices, and take the total number of relations
to be \(M_P=P(P-1)/2=KN(KN-1)/2\). The phase cycle is closed by
\(U^K=I\). In this construction, even for \(N=1\), as long as
\(K\ge 3\), multiple relations and cycles can be defined among the
distinct phase states of a single object. This paper completely analyzes
the minimal nontrivial exactly closed case \(N=1, K=4\). Here
\(U=\exp(2\pi i/4)=i\) is taken as the primitive 4-th root, and the
states are defined by \(S_k=(U^k,U^{mk})\). We compare \(m=2\), the
minimal nontrivial even harmonic, with \(m=3\), the minimal nontrivial
odd harmonic.

The results are stated in three strictly separated layers. First, the
existence of the complete graph \(K_4\), the state cycle, the orders,
and the squared sums of the two wave components are mathematically exact
and require no physical identification. Second, when the Euclidean
distance \(d_{ij}\) on \(\mathbb C^2\) is \textbf{defined as an
analytical readout}, the distance multisets of the 6 edges differ
between \(m=2\) and \(m=3\). This is exact once defined, but identifying
\(d_{ij}\) with real spatial distance is not a claim of this paper.
Third, identifying the \(N=1\) system with ``one photon'', reading the
cycle of \(k\) as physical past, present, and future, and reading the
network distance as physical length are all \textbf{physical hypotheses}
to be tested in the future; they are not preconditions for the
complete-network analysis of this paper.

For \(K=4\), the minimal order of the harmonic component degenerates to
2 when \(m=2\), while both the base wave and the harmonic retain minimal
order 4 when \(m=3\). Furthermore, defining the total squared sum of the
two wave components over the 4 states as

\[
Q_{\mathrm{loop}}(m)
=
\sum_{k=0}^{3}\left[(U^k)^2+(U^{mk})^2\right]
\]

we obtain exactly

\[
Q_{\mathrm{loop}}(m)
=
\begin{cases}
4,&m\ \mathrm{even},\\
0,&m\ \mathrm{odd}
\end{cases}
\]

Thus, in the minimal comparison, \(m=2\) does not close to squared-sum
zero, whereas \(m=3\) closes to squared-sum zero exactly. This paper
derives this difference completely --- not as a numerical approximation
but as geometric series over a finite cyclic group --- and includes all
states, all 6 relations, all visualizations, dynamic HTML, MP4, CSV, and
reproduction programs in the same folder.

\textbf{Keywords:} complete graph, discrete phase, finite cyclic group,
\(U^K=I\), relational waves, time as vertices, zero square-sum closure,
harmonics, anonymity, self-consistency

\begin{center}\rule{0.5\linewidth}{0.5pt}\end{center}

\section{1. Fixing the Scope of Claims
First}\label{fixing-the-scope-of-claims-first}

This paper deliberately separates physical identification from
mathematical network analysis. This is not a stylistic caution; it is
the logical structure of the paper itself.

{\def\LTcaptype{none} % do not increment counter
\begin{longtable}[]{@{}
  >{\raggedright\arraybackslash}p{(\linewidth - 4\tabcolsep) * \real{0.3333}}
  >{\raggedright\arraybackslash}p{(\linewidth - 4\tabcolsep) * \real{0.3333}}
  >{\raggedright\arraybackslash}p{(\linewidth - 4\tabcolsep) * \real{0.3333}}@{}}
\toprule\noalign{}
\begin{minipage}[b]{\linewidth}\raggedright
Layer
\end{minipage} & \begin{minipage}[b]{\linewidth}\raggedright
Content
\end{minipage} & \begin{minipage}[b]{\linewidth}\raggedright
Status in this paper
\end{minipage} \\
\midrule\noalign{}
\endhead
\bottomrule\noalign{}
\endlastfoot
A & Treating \(P=KN\) vertices with all relations, the 6 relations of
\(K_4\), \(U^K=I\), the state cycle, orders, squared sums of the two
wave components & \textbf{Exact mathematical/algebraic claims} \\
B & Numerical edge lengths of the 6 edges via the Euclidean distance
\(d_{ij}\) on \(\mathbb C^2\) & \textbf{Explicitly defined analytical
readout}. Exact once defined, but not physical distance \\
C & 3-dimensional configuration via MDS & \textbf{Visualization of the
distance matrix}. No physical meaning is assigned to the coordinate
values themselves \\
D & Regarding \(N=1\) as one photon, and the base wave plus one harmonic
as its minimal state & \textbf{Physical hypothesis} \\
E & Reading the phase state \(k\) as physical past/present/future, and
reading network distance as spacetime distance & \textbf{Physical
readout hypotheses to be tested in the future} \\
\end{longtable}
}

Therefore, even if D or E is refuted in the future, the complete-network
analysis of layer A and the results on the finite cyclic group are
unaffected. Likewise, even if the distance readout of layer B turns out
to be unusable as physical length, the complete relational structure of
\(K_4\), the orders, and the squared-sum results remain.

\begin{center}\rule{0.5\linewidth}{0.5pt}\end{center}

\section{2. Motivation: What Was Observed in the Preceding Inflation
System}\label{motivation-what-was-observed-in-the-preceding-inflation-system}

\subsection{2.1 Basic construction of the preceding
system}\label{basic-construction-of-the-preceding-system}

In the preceding study {[}K1{]}, \(N\) vertices were nominally placed
``at an instant'', and a complex relational wave was assigned to every
vertex pair of the complete graph. The number of relations is

\[
M
=
\binom{N}{2}
=
\frac{N(N-1)}{2}
\]

For example, with \(N=40\),

\[
M
=
\frac{40\times39}{2}
=
780
\]

relational waves exist.

In the preceding experiments, when these relational waves are observed
in the complex plane, highly symmetric initial states concentrate near
the 90-degree series and form a cross-shaped set of narrow bands.
Because this is easily misunderstood from text alone, we reproduce the
measured figure for \(N=40\) as is.

\subsubsection{\texorpdfstring{Figure 1: Initial state at
\(N=40, M=780\): cross-shaped bands near the 90-degree
series}{Figure 1: Initial state at N=40, M=780: cross-shaped bands near the 90-degree series}}\label{figure-1-initial-state-at-n40-m780-cross-shaped-bands-near-the-90-degree-series}

\begin{figure}
\centering
\pandocbounded{\includegraphics[keepaspectratio,alt={N=40 step0 cross phase distribution}]{figs/inflation_N40_phase_step0_cross_en_20260915.png}}
\caption{N=40 step0 cross phase distribution}
\end{figure}

In Figure 1, the relational waves are not placed uniformly on a phase
circle; they are strongly constrained near four directions about 90
degrees apart, each forming a band with a small amplitude width. Thus
the statement ``the initial state is cross-shaped'' is not a schematic
metaphor but an observation obtained by directly plotting the set of
\(N=40, M=780\) relational waves in the complex plane.

\subsection{2.2 Quasi-stable ring after
iteration}\label{quasi-stable-ring-after-iteration}

In the preceding system, iterating the interaction causes a rapid growth
of the transverse component, after which the amplitudes become nearly
uniform and the phase constraint largely dissolves. The terminal
distribution is, however, not a perfectly equiangular configuration; it
retains an anisotropy corresponding to the original 90-degree series.

\subsubsection{\texorpdfstring{Figure 2: Terminal state at
\(N=40, M=780\): radial
ring}{Figure 2: Terminal state at N=40, M=780: radial ring}}\label{figure-2-terminal-state-at-n40-m780-radial-ring}

\begin{figure}
\centering
\pandocbounded{\includegraphics[keepaspectratio,alt={N=40 step500 ring phase distribution}]{figs/inflation_N40_phase_step500_ring_en_20260915.png}}
\caption{N=40 step500 ring phase distribution}
\end{figure}

Figure 2 is a terminal example at step 500. The distribution has
transitioned from the initial four narrow bands to a ring of nearly
equal radius, but the angular density is not perfectly uniform. This
paper takes this phenomenon as the starting problem: can the totality of
state relations, including the time direction, be described from the
outset as a single closed network?

\subsection{2.3 What ``inflation-like'' means in this
paper}\label{what-inflation-like-means-in-this-paper}

When this paper uses the words ``inflation'' or ``inflation-like'', it
\textbf{does not presuppose identity with cosmological inflation}. It
refers to the observed behavior in the preceding relational-wave system
in which, over a short interval starting from a small transverse
component, the transverse order parameter grows rapidly and the system
moves from a symmetric unstable state to a degenerate quasi-stable
state.

To make the meaning of this terminology clear, we reproduce two figures
from the preceding experiment.

\subsubsection{Figure 3: Mexican-hat SSB figure reconstructed from
measured order-parameter
dynamics}\label{figure-3-mexican-hat-ssb-figure-reconstructed-from-measured-order-parameter-dynamics}

\begin{figure}
\centering
\pandocbounded{\includegraphics[keepaspectratio,alt={Mexican-hat SSB reconstructed from measured dynamics}]{figs/inflation_mexican_hat_SSB_reconstruction_20260915.png}}
\caption{Mexican-hat SSB reconstructed from measured dynamics}
\end{figure}

This figure was not drawn by assuming a schematic potential in advance.
It is an auxiliary visualization reconstructed by mapping the evolution
of the measured order parameter

\[
r=\sqrt{H_\perp/H}
\]

onto a Landau-type effective potential

\[
V(r)
=
-\frac12\lambda_0 r^2
+
\frac{\lambda_0}{4r_{\mathrm{vac}}^2}r^4
\]

with coefficients set from measured values as

\[
\lambda_0=\ln\rho_0,
\qquad
r_{\mathrm{vac}}
=
\sqrt{(H_\perp/H)_{\mathrm{sat}}}
\]

Here \(H=\langle z,z\rangle=\sum_e|z_e|^2\) is the total squared norm of
the relational-wave state, and \(H_\perp=\|\Pi_\perp z\|^2\) is the
squared norm of the component in the orthogonal complement of the
initial floor plane \(\Pi_0\). Thus \(H_\perp/H\) is a dimensionless
readout quantity representing the relative departure from the initial
floor plane. Further, \(J_0\) is the Jacobian in the co-rotating
coordinates that fix the initial relative equilibrium,
\(\rho_0=\rho(J_0)\) is its spectral radius, and
\((H_\perp/H)_{\mathrm{sat}}\) is the saturation value of \(H_\perp/H\)
measured in the terminal quasi-stable phase.

Therefore this \(V(r)\) is an \textbf{effective description (ansatz)
whose coefficients are identified from measured quantities}; the
potential shape itself is not derived from first principles in this
paper. Nor is this effective potential a precondition of the \(K_4\)
complete-network analysis. What matters here is to show visually the
preceding observation of the transition from the central high-symmetry
state to the terminal ring-shaped quasi-stable set. This paper
introduces no new dynamical law from this figure.

\subsubsection{\texorpdfstring{Figure 4: 3D evolution display of the
\(N=40\) relational
waves}{Figure 4: 3D evolution display of the N=40 relational waves}}\label{figure-4-3d-evolution-display-of-the-n40-relational-waves}

\begin{figure}
\centering
\pandocbounded{\includegraphics[keepaspectratio,alt={N=40 relation-wave 3D expansion}]{figs/inflation_N40_relation_wave_3D_step500_en_20260915.png}}
\caption{N=40 relation-wave 3D expansion}
\end{figure}

Figure 4 is an example displaying, in three dimensions, the trajectories
of many relational waves from the initial highly symmetric constrained
state to the terminal ring. This, too, is not an input to the \(K_4\)
analysis of this paper. It is used \textbf{only as the preceding
observation that motivated this paper}.

\begin{center}\rule{0.5\linewidth}{0.5pt}\end{center}

\section{3. Anonymity Concerns Remaining in the Preceding
System}\label{anonymity-concerns-remaining-in-the-preceding-system}

By \textbf{anonymity} this paper means assigning to none of the states
and variables constituting the system any a priori kind (for example,
``spatial vertex'' versus ``time label'') or any external role (for
example, a special variable that provides the order in which the other
states are updated). Accordingly, an ``anonymity concern'' in this paper
means a structural check of whether some variable or vertex set is being
treated specially from outside the computational procedure.

\subsection{\texorpdfstring{3.1 Special treatment of the external update
variable
\(\tau\)}{3.1 Special treatment of the external update variable \textbackslash tau}}\label{special-treatment-of-the-external-update-variable-tau}

In the preceding experiments, the interaction was iterated along a
discrete update variable \(\tau\), typically with

\[
\Delta\tau
=
\frac{2\pi}{N}
\]

Here, this paper does not assume that \(\tau\) is identified with
physical time. Even so, as long as the computational procedure updates
one way,

\[
\text{state at }\tau
\longrightarrow
\text{state at }\tau+\Delta\tau
\]

\(\tau\) alone is treated specially as ``the external order used to
update the other variables''.

If anonymity is to be enforced thoroughly, the question remains:

\[
\boxed{
\text{why must }\tau\text{ alone be an external label — neither a vertex nor a relation?}
}
\]

\subsection{\texorpdfstring{3.2 The assumption of ``\(N\) vertices at an
instant''}{3.2 The assumption of ``N vertices at an instant''}}\label{the-assumption-of-n-vertices-at-an-instant}

Similarly, when we posit

\[
N\text{ vertices}
\quad\Longrightarrow\quad
M=\frac{N(N-1)}{2}\text{ relations}
\]

we classify the \(N\) vertices in advance as ``existing at the same
instant''. Seen from a higher-level anonymous relational system, this
classification may also be artificial.

Therefore, in this paper, ``vertices arranged in space'' and ``states
arranged along the time direction'' are not treated as different kinds
from the outset.

\begin{center}\rule{0.5\linewidth}{0.5pt}\end{center}

\section{4. Turning Time-Direction States into
Vertices}\label{turning-time-direction-states-into-vertices}

\subsection{\texorpdfstring{4.1 Total number of vertices
\(P=KN\)}{4.1 Total number of vertices P=KN}}\label{total-number-of-vertices-pkn}

Nominally, suppose one phase section carries \(N\) states, and the
closed phase sequence carries \(K\) state labels.

Conventionally, one treats

\[
\{1,2,\ldots,N\}
\]

as ``vertices'' and

\[
\{0,1,\ldots,K-1\}
\]

separately as ``time or update labels''.

This paper removes this distinction and counts the pairs

\[
(k,a),
\qquad
k\in\{0,1,\ldots,K-1\},
\qquad
a\in\{1,2,\ldots,N\}
\]

all as vertices of the same kind.

Therefore the total number of vertices is

\[
P
=
K\times N
=
KN
\]

No assumption ``\(k\) is physical time'' enters here. \(k\) is a state
label on a finite closed cycle; whether it is later read as physical
time is a separate question.

\subsection{4.2 Derivation of the total number of
relations}\label{derivation-of-the-total-number-of-relations}

Keeping every distinct pair of the \(P\) vertices as a relation, the
number of edges of the complete graph \(K_P\) is

\[
M_P
=
\binom{P}{2}
\]

From the definition of the binomial coefficient,

\[
\binom{P}{2}
=
\frac{P!}{2!(P-2)!}
\]

and since

\[
P!
=
P(P-1)(P-2)!
\]

we have

\[
\binom{P}{2}
=
\frac{P(P-1)(P-2)!}{2(P-2)!}
=
\frac{P(P-1)}{2}.
\]

Substituting \(P=KN\),

\[
\boxed{
M_P
=
\frac{KN(KN-1)}{2}
}
\]

This formula assumes no physical meaning of time. It is the
combinatorial count obtained when ``\(N\) objects with \(K\) distinct
state layers'' are completely related as \(KN\) vertices without
distinction.

\subsection{4.3 Why we do not directly compute huge
systems}\label{why-we-do-not-directly-compute-huge-systems}

Conventional numerical experiments reach at most \(N\sim40\). Decomposed
as \(KN\), even \(N=6, K=7\) already gives a comparable number of
vertices.

On the other hand, if one considers \(N\sim10^{60}\) as a rough physical
state count and also \(K\sim10^{60}\) in the phase direction, then

\[
P=KN\sim10^{120}
\]

and the number of complete relations is

\[
M_P
=
\frac{P(P-1)}{2}
\sim
\frac{10^{240}}{2}
\]

that is, of order \(10^{240}\). These numbers are not experimental
values of this paper; they are order estimates showing that direct
enumeration of huge systems is impossible in principle.

Therefore the strategy of this paper is not to coarsely approximate a
huge system, but to

\[
\boxed{
\text{solve completely the minimal exactly closed system, keeping the same closure axiom and complete relational structure}
}
\]

\begin{center}\rule{0.5\linewidth}{0.5pt}\end{center}

\section{\texorpdfstring{5. Why Can We Go Down to
\(N=1\)?}{5. Why Can We Go Down to N=1?}}\label{why-can-we-go-down-to-n1}

\subsection{5.1 Looking only at a single
section}\label{looking-only-at-a-single-section}

Considering the ordinary complete graph alone:

\subsubsection{\texorpdfstring{\(N=1\)}{N=1}}\label{n1}

\[
M
=
\frac{1(1-1)}{2}
=0.
\]

There is no relation.

\subsubsection{\texorpdfstring{\(N=2\)}{N=2}}\label{n2}

\[
M
=
\frac{2(2-1)}{2}
=1.
\]

There is only one relation, so relations cannot be compared with each
other.

\subsubsection{\texorpdfstring{\(N=3\)}{N=3}}\label{n3}

\[
M
=
\frac{3(3-1)}{2}
=3.
\]

For the first time, three relations

\[
r_{12},\quad r_{23},\quad r_{31}
\]

can be mutually compared, and the cycle

\[
1\to2\to3\to1
\]

can be defined.

Therefore, counting only ``vertices at the same nominal instant'',
\(N=3\) is the minimal nontrivial cycle.

\subsection{\texorpdfstring{5.2 With state-direction vertices, a cycle
exists already at
\(N=1, K\ge3\)}{5.2 With state-direction vertices, a cycle exists already at N=1, K\textbackslash ge3}}\label{with-state-direction-vertices-a-cycle-exists-already-at-n1-kge3}

However, if the \(K\) states in the state direction are all vertices,
then even for \(N=1\) the total number of vertices is

\[
P=K.
\]

The number of relations is then

\[
M_P
=
\frac{K(K-1)}{2}.
\]

For \(K=3\),

\[
M_P
=
3,
\]

for \(K=4\),

\[
M_P
=
6.
\]

That is, even a single object, if its distinct phase states are treated
as equal vertices, carries multiple relations among its own distinct
states and can form a cycle.

The initial state of the preceding system was constrained, as shown in
Figure 1, to four directional bands about 90 degrees apart in the
complex plane. As the minimal phase cycle corresponding to this
\textbf{observed four-direction structure}, we adopt

\[
\boxed{N=1,\quad K=4}
\]

On top of this, we compare the minimal nontrivial even harmonic \(m=2\)
and odd harmonic \(m=3\). Thus ``\(K=4\) because of 90 degrees'' is not
placed circularly from the definition alone; it is a comparison setting
that maps the four-direction structure observed in the preceding
experiment into the minimal exactly closed system of this paper.

This is not a coarse approximation of a huge system. It is a
\textbf{finite complete-relation closed system that keeps all 4 vertices
and all 6 relations without omitting a single one}.

\begin{center}\rule{0.5\linewidth}{0.5pt}\end{center}

\section{6. Physical Correspondences Are Kept Separate as
Hypotheses}\label{physical-correspondences-are-kept-separate-as-hypotheses}

\subsection{6.1 ``One photon'' is not a precondition of the mathematical
analysis}\label{one-photon-is-not-a-precondition-of-the-mathematical-analysis}

As a physical candidate for \(N=1\), this paper envisions ``one
photon''. However, this identification is unnecessary for the
complete-network analysis to hold.

We therefore clearly distinguish:

\[
\boxed{
\text{analyzing the }N=1\text{ finite closed system}
}
\]

is the mathematical setting of this paper, whereas

\[
\boxed{\text{reading that }N=1\text{ as one photon}}
\]

is a physical hypothesis.

\subsection{6.2 Motivation: a single complex wave is insufficient for
localization
readout}\label{motivation-a-single-complex-wave-is-insufficient-for-localization-readout}

Let a single complex wave be

\[
z
=
a+ib
\]

In polar form,

\[
z
=
A e^{i\theta}
\]

If, as a single monochromatic wave, \(A\) is constant, its phase exists
over the entire closed cycle, and one wave alone carries little
information for distinguishing a localized state position.

This paper therefore uses, as the minimal additional relation, ``a base
wave plus one integer harmonic'':

\[
z_1(\theta)
=
e^{i\theta},
\]

\[
z_m(\theta)
=
e^{im\theta}.
\]

This two-wave system shares the same phase variable \(\theta\). No
independent second clock is introduced.

Moreover, in this paper the amplitudes are defined to be 1, and this is
not a normalization operation performed at each step. That is,

\[
|z_1|=|z_m|=1
\]

is \textbf{the definition of the state family itself}; no operation
readjusting the norm during time evolution is introduced.

\begin{center}\rule{0.5\linewidth}{0.5pt}\end{center}

\section{\texorpdfstring{7. Definition of the Finite Cycle:
\(U^K=I\)}{7. Definition of the Finite Cycle: U\^{}K=I}}\label{definition-of-the-finite-cycle-uki}

\subsection{\texorpdfstring{7.1 Choosing a primitive \(K\)-th
root}{7.1 Choosing a primitive K-th root}}\label{choosing-a-primitive-k-th-root}

The condition

\[
U^K=I
\]

alone does not guarantee that the minimal order of \(U\) is \(K\). For
example, \(U=1\) also satisfies \(U^K=1\) but closes in one step.

Therefore, to obtain \(K\) distinct base phases, this paper adopts the
primitive \(K\)-th root

\[
U
=
\exp\left(\frac{2\pi i}{K}\right)
\]

Then

\[
U^K
=
\exp\left(\frac{2\pi i}{K}K\right)
=
\exp(2\pi i)
=
1.
\]

Also, for \(0<q<K\), in general \(U^q\ne1\), so the minimal order of the
base generator is \(K\).

\subsection{\texorpdfstring{7.2 In this paper,
\(K=4\)}{7.2 In this paper, K=4}}\label{in-this-paper-k4}

\[
U
=
\exp\left(\frac{2\pi i}{4}\right)
=
\exp\left(\frac{\pi i}{2}\right)
=i.
\]

Therefore

\[
U^0=1,
\]

\[
U^1=i,
\]

\[
U^2=-1,
\]

\[
U^3=-i,
\]

\[
U^4=1.
\]

The four discrete phase states are labeled

\[
k=0,1,2,3
\]

\subsection{7.3 Definition of the two-wave
states}\label{definition-of-the-two-wave-states}

For an integer harmonic \(m\), define the state \(S_k\) by

\[
\boxed{
S_k
=
\left(z_1(k),z_m(k)\right)
=
\left(U^k,U^{mk}\right)
}
\]

From this definition, the next state is

\[
S_{k+1}
=
\left(U^{k+1},U^{m(k+1)}\right).
\]

Using the law of exponents,

\[
U^{k+1}
=
U\,U^k,
\]

\[
U^{m(k+1)}
=
U^{mk+m}
=
U^m U^{mk}.
\]

Therefore

\[
\boxed{
S_{k+1}
=
\left(U z_1(k),U^m z_m(k)\right)
}
\]

Furthermore, after \(K\) steps,

\[
S_{k+K}
=
\left(U^{k+K},U^{m(k+K)}\right).
\]

The first component is

\[
U^{k+K}
=
U^k U^K
=
U^k.
\]

The second component is

\[
U^{m(k+K)}
=
U^{mk}U^{mK}
=
U^{mk}(U^K)^m
=
U^{mk}.
\]

Hence

\[
\boxed{
S_{k+K}=S_k
}
\]

Here the fact that ``all states form a cycle'' is derived algebraically
from the definition.

\begin{center}\rule{0.5\linewidth}{0.5pt}\end{center}

\section{\texorpdfstring{8. Minimal Comparison: Even Harmonic \(m=2\)
versus Odd Harmonic
\(m=3\)}{8. Minimal Comparison: Even Harmonic m=2 versus Odd Harmonic m=3}}\label{minimal-comparison-even-harmonic-m2-versus-odd-harmonic-m3}

\subsection{\texorpdfstring{8.1 Even harmonic
\(m=2\)}{8.1 Even harmonic m=2}}\label{even-harmonic-m2}

The definition is

\[
S_k
=
(U^k,U^{2k}).
\]

\subsubsection{\texorpdfstring{\(k=0\)}{k=0}}\label{k0}

\[
S_0
=
(U^0,U^0)
=
(1,1).
\]

\subsubsection{\texorpdfstring{\(k=1\)}{k=1}}\label{k1}

\[
S_1
=
(U,U^2)
=
(i,-1).
\]

\subsubsection{\texorpdfstring{\(k=2\)}{k=2}}\label{k2}

\[
S_2
=
(U^2,U^4)
=
(-1,1).
\]

\subsubsection{\texorpdfstring{\(k=3\)}{k=3}}\label{k3}

\[
S_3
=
(U^3,U^6).
\]

Here

\[
U^6
=
U^{4+2}
=
U^4U^2
=
1\times(-1)
=
-1
\]

so

\[
S_3
=
(-i,-1).
\]

Therefore

\[
\boxed{
(1,1)
\to
(i,-1)
\to
(-1,1)
\to
(-i,-1)
\to
(1,1)
}
\]

Looking at the harmonic component alone,

\[
1,-1,1,-1
\]

it has degenerated to 2 states instead of 4.

\subsection{\texorpdfstring{8.2 Odd harmonic
\(m=3\)}{8.2 Odd harmonic m=3}}\label{odd-harmonic-m3}

The definition is

\[
S_k
=
(U^k,U^{3k}).
\]

\subsubsection{\texorpdfstring{\(k=0\)}{k=0}}\label{k0-1}

\[
S_0
=
(1,1).
\]

\subsubsection{\texorpdfstring{\(k=1\)}{k=1}}\label{k1-1}

\[
S_1
=
(U,U^3)
=
(i,-i).
\]

\subsubsection{\texorpdfstring{\(k=2\)}{k=2}}\label{k2-1}

\[
S_2
=
(U^2,U^6).
\]

From

\[
U^6
=
U^{4+2}
=
U^2
=
-1
\]

we get

\[
S_2
=
(-1,-1).
\]

\subsubsection{\texorpdfstring{\(k=3\)}{k=3}}\label{k3-1}

\[
S_3
=
(U^3,U^9).
\]

Since

\[
U^9
=
U^{8+1}
=
(U^4)^2U
=
i
\]

we have

\[
S_3
=
(-i,i).
\]

Therefore

\[
\boxed{
(1,1)
\to
(i,-i)
\to
(-1,-1)
\to
(-i,i)
\to
(1,1)
}
\]

In this case, the harmonic component also cycles through all four
states:

\[
1,-i,-1,i
\]

\begin{center}\rule{0.5\linewidth}{0.5pt}\end{center}

\section{9. Algebraic Self-Consistency as a Finite Closed
System}\label{algebraic-self-consistency-as-a-finite-closed-system}

To avoid confusing the term ``self-consistent'' with the full
nonlinear-interaction fixed-point condition of the preceding inflation
system, this paper restricts its meaning.

What is checked here is \textbf{algebraic consistency as a finite closed
system}, namely:

\begin{enumerate}
\def\labelenumi{\arabic{enumi}.}
\tightlist
\item
  As defined, \(|z_1|=|z_m|=1\).
\item
  At every state, \(z_m=z_1^m\) holds.
\item
  The transition rule \(S_{k+1}=(Uz_1,U^m z_m)\) holds.
\item
  Forward and reverse give the same complex state for the same \(k\).
\item
  After \(K\) steps the system returns exactly to the initial state.
\end{enumerate}

\subsection{9.1 Amplitude condition}\label{amplitude-condition}

Since \(U\) lies on the unit circle,

\[
|U|=1.
\]

Therefore, for any integer \(q\),

\[
|U^q|
=
|U|^q
=
1.
\]

Hence

\[
|z_1(k)|
=
|U^k|
=
1,
\]

\[
|z_m(k)|
=
|U^{mk}|
=
1.
\]

\subsection{9.2 Harmonic relation}\label{harmonic-relation}

The first component is

\[
z_1(k)=U^k.
\]

Its \(m\)-th power is

\[
[z_1(k)]^m
=
(U^k)^m
=
U^{mk}
=
z_m(k).
\]

Therefore

\[
\boxed{z_m=z_1^m}
\]

holds identically at every state.

\subsection{9.3 Forward / reverse}\label{forward-reverse}

Let forward be

\[
0\to1\to2\to3
\]

and reverse be

\[
0\to3\to2\to1
\]

Since the state \(S_k\) is uniquely defined as a function of \(k\), not
of the ordering, forward and reverse give the same \(S_k\) for the same
\(k\).

This is not a proof that ``physical time is reversible''. It is a check
that \textbf{enumerating the same finite cycle in two directions
produces no contradiction in the state definition}.

\subsection{9.4 Cross-check against numerical
data}\label{cross-check-against-numerical-data}

In the raw data generated with exact-\(i\), the maximum residual of the
above algebraic conditions is 0 for both \(m=2\) and \(m=3\).

\begin{itemize}
\tightlist
\item
  \href{reproducibility_raw_readout_20260915/raw_data/complex_states_animation_exact_i.csv}{exact-\(i\)
  complex state data}
\item
  \href{reproducibility_raw_readout_20260915/raw_data/complex_states_static_exp.csv}{static
  recomputation via \(\exp(i\theta)\)}
\item
  \href{reproducibility_raw_readout_20260915/self_consistency_square_sum_U4_20260915/K4_self_consistency_square_sum_by_state_20260915.csv}{per-state
  check CSV for self-consistency and squared sums}
\item
  \href{reproducibility_raw_readout_20260915/self_consistency_square_sum_U4_20260915/K4_self_consistency_square_sum_summary_20260915.csv}{summary
  CSV}
\item
  \href{reproducibility_raw_readout_20260915/self_consistency_square_sum_U4_20260915/K4_self_consistency_square_sum_analysis_20260915.md}{analysis
  details}
\end{itemize}

The maximum difference from floating-point \(\exp(i\theta)\) generation
is about \(4.0\times10^{-16}\) for \(m=2\) and about
\(5.5\times10^{-16}\) for \(m=3\), in agreement with the exact-\(i\)
discrete states within rounding error.

\begin{center}\rule{0.5\linewidth}{0.5pt}\end{center}

\section{\texorpdfstring{10. The Complete Network
\(K_4\)}{10. The Complete Network K\_4}}\label{the-complete-network-k_4}

\subsection{10.1 Not a single one of the 6 relations is
omitted}\label{not-a-single-one-of-the-6-relations-is-omitted}

Since \(K=4\), there are 4 vertices, and the number of edges of the
complete graph is

\[
M
=
\frac{4(4-1)}{2}
=
6.
\]

All vertex pairs are

\[
(0,1),
(0,2),
(0,3),
(1,2),
(1,3),
(2,3)
\]

The important point is that the existence of the complete graph itself
requires no definition of ``length''. The combinatorial structure of
keeping all 6 relations comes first.

\subsection{\texorpdfstring{10.2 The distance on \(\mathbb C^2\) is
defined as an analytical
readout}{10.2 The distance on \textbackslash mathbb C\^{}2 is defined as an analytical readout}}\label{the-distance-on-mathbb-c2-is-defined-as-an-analytical-readout}

For states

\[
S_i=(U^i,U^{mi}),
\]

\[
S_j=(U^j,U^{mj})
\]

take the difference

\[
S_i-S_j
=
(U^i-U^j,\;U^{mi}-U^{mj})
\]

Defining, \textbf{as an analytical readout}, the standard Euclidean norm
of this two-component complex vector,

\[
\boxed{
d_{ij}
=
\sqrt{
|U^i-U^j|^2
+
|U^{mi}-U^{mj}|^2
}
}
\]

This paper gives no grounds for identifying \(d_{ij}\) with real spatial
distance, spacetime interval, or interparticle distance. The distance
geometry below is therefore \textbf{state-space geometry} based on the
definition on \(\mathbb C^2\).

\subsection{10.3 Complete derivation of the distance
formula}\label{complete-derivation-of-the-distance-formula}

In general, for unit complex numbers \(e^{i\alpha},e^{i\beta}\),

\[
|e^{i\alpha}-e^{i\beta}|^2
=
(e^{i\alpha}-e^{i\beta})(e^{-i\alpha}-e^{-i\beta}).
\]

Expanding,

\[
=
1-e^{i(\alpha-\beta)}-e^{-i(\alpha-\beta)}+1.
\]

By Euler's formula,

\[
e^{ix}+e^{-ix}=2\cos x
\]

so

\[
|e^{i\alpha}-e^{i\beta}|^2
=
2-2\cos(\alpha-\beta).
\]

Together with

\[
1-\cos x
=
2\sin^2\frac{x}{2}
\]

we get

\[
|e^{i\alpha}-e^{i\beta}|^2
=
4\sin^2\frac{\alpha-\beta}{2}.
\]

With \(U=e^{2\pi i/4}\) and \(\Delta=i-j\), the first component gives

\[
|U^i-U^j|^2
=
4\sin^2\left(\frac{\pi\Delta}{4}\right).
\]

The second component gives

\[
|U^{mi}-U^{mj}|^2
=
4\sin^2\left(\frac{\pi m\Delta}{4}\right).
\]

Therefore

\[
\boxed{
d_{ij}^2
=
4\sin^2\left(\frac{\pi\Delta}{4}\right)
+
4\sin^2\left(\frac{\pi m\Delta}{4}\right)
}
\]

\subsection{\texorpdfstring{10.4 All 6 edges for
\(m=2\)}{10.4 All 6 edges for m=2}}\label{all-6-edges-for-m2}

For \(\Delta=1\),

\[
4\sin^2\frac{\pi}{4}=2,
\]

\[
4\sin^2\frac{2\pi}{4}
=
4\sin^2\frac{\pi}{2}
=4.
\]

Therefore

\[
d^2=2+4=6,
\]

\[
d=\sqrt6.
\]

For \(\Delta=2\),

\[
4\sin^2\frac{2\pi}{4}
=4,
\]

\[
4\sin^2\frac{4\pi}{4}
=
4\sin^2\pi
=0.
\]

Therefore

\[
d^2=4,
\]

\[
d=2.
\]

\(\Delta=3\) is symmetric to \(\Delta=1\), so \(d=\sqrt6\).

All 6 edges are as follows.

{\def\LTcaptype{none} % do not increment counter
\begin{longtable}[]{@{}lrr@{}}
\toprule\noalign{}
Edge & \(\min(\lvert\Delta\rvert,\,4-\lvert\Delta\rvert)\) &
\(d_{ij}\) \\
\midrule\noalign{}
\endhead
\bottomrule\noalign{}
\endlastfoot
\((0,1)\) & 1 & \(\sqrt6\) \\
\((0,2)\) & 2 & \(2\) \\
\((0,3)\) & 1 & \(\sqrt6\) \\
\((1,2)\) & 1 & \(\sqrt6\) \\
\((1,3)\) & 2 & \(2\) \\
\((2,3)\) & 1 & \(\sqrt6\) \\
\end{longtable}
}

Hence the distance multiset is

\[
\boxed{
\{2,2,\sqrt6,\sqrt6,\sqrt6,\sqrt6\}
}
\]

\subsubsection{\texorpdfstring{Figure 5: Complete state network for
\(m=2\)}{Figure 5: Complete state network for m=2}}\label{figure-5-complete-state-network-for-m2}

\begin{figure}
\centering
\pandocbounded{\includegraphics[keepaspectratio,alt={K4 complete network m2}]{figs/K4_complete_network_m2.png}}
\caption{K4 complete network m2}
\end{figure}

\href{K4_complete_network_m2.svg}{SVG version} /
\href{edges_m2.csv}{edge data CSV}

\subsection{\texorpdfstring{10.5 All 6 edges for
\(m=3\)}{10.5 All 6 edges for m=3}}\label{all-6-edges-for-m3}

For \(\Delta=1\),

\[
4\sin^2\frac{\pi}{4}=2,
\]

\[
4\sin^2\frac{3\pi}{4}=2.
\]

Therefore

\[
d^2=4,
\]

\[
d=2.
\]

For \(\Delta=2\),

\[
4\sin^2\frac{2\pi}{4}=4,
\]

\[
4\sin^2\frac{6\pi}{4}
=
4\sin^2\frac{3\pi}{2}
=4.
\]

Therefore

\[
d^2=8,
\]

\[
d=2\sqrt2.
\]

\(\Delta=3\) is symmetric to \(\Delta=1\), so \(d=2\).

All 6 edges are as follows.

{\def\LTcaptype{none} % do not increment counter
\begin{longtable}[]{@{}lrr@{}}
\toprule\noalign{}
Edge & \(\min(\lvert\Delta\rvert,\,4-\lvert\Delta\rvert)\) &
\(d_{ij}\) \\
\midrule\noalign{}
\endhead
\bottomrule\noalign{}
\endlastfoot
\((0,1)\) & 1 & \(2\) \\
\((0,2)\) & 2 & \(2\sqrt2\) \\
\((0,3)\) & 1 & \(2\) \\
\((1,2)\) & 1 & \(2\) \\
\((1,3)\) & 2 & \(2\sqrt2\) \\
\((2,3)\) & 1 & \(2\) \\
\end{longtable}
}

Hence the distance multiset is

\[
\boxed{
\{2,2,2,2,2\sqrt2,2\sqrt2\}
}
\]

\subsubsection{\texorpdfstring{Figure 6: Complete state network for
\(m=3\)}{Figure 6: Complete state network for m=3}}\label{figure-6-complete-state-network-for-m3}

\begin{figure}
\centering
\pandocbounded{\includegraphics[keepaspectratio,alt={K4 complete network m3}]{figs/K4_complete_network_m3.png}}
\caption{K4 complete network m3}
\end{figure}

\href{K4_complete_network_m3.svg}{SVG version} /
\href{edges_m3.csv}{edge data CSV} / \href{all_edges_m2_m3.csv}{combined
\(m=2,3\) edge data}

\subsection{10.6 What the MDS display means --- and what it does not
claim}\label{what-the-mds-display-means-and-what-it-does-not-claim}

The 3-dimensional coordinates in Figures 5 and 6 are display coordinates
reconstructed by classical MDS from the distance matrix of the 6 values
\(d_{ij}\).

Therefore it is essential that

\[
\boxed{
\text{the MDS-1, MDS-2, MDS-3 axes of the figures are not identified with physical }x,y,z
}
\]

What this paper claims to be meaningful is not the 3D coordinates
themselves but the fact that \textbf{for the defined \(\mathbb C^2\)
distance, the distance multisets of the 6 edges differ between \(m=2\)
and \(m=3\)}.

\begin{center}\rule{0.5\linewidth}{0.5pt}\end{center}

\section{11. Waveform Displays: Continuous Curves Are Readout Figures;
the States Themselves Are
Discrete}\label{waveform-displays-continuous-curves-are-readout-figures-the-states-themselves-are-discrete}

The state network consists of only the 4 points

\[
k=0,1,2,3
\]

On the other hand, to understand the waveforms visually, the same base
and harmonic can be drawn on the continuous phase variable \(\theta\) as

\[
\psi_1(\theta)=\cos\theta,
\]

\[
\psi_m(\theta)=\cos(m\theta)
\]

Drawing continuous curves here is not a new assumption that the actual
state transitions are continuous. The discrete states of this paper
remain the 4 states sampled at

\[
\theta_k
=
\frac{2\pi k}{4}
=
\frac{k\pi}{2}
\]

\subsubsection{\texorpdfstring{Figure 7: Base wave and even harmonic
\(m=2\)}{Figure 7: Base wave and even harmonic m=2}}\label{figure-7-base-wave-and-even-harmonic-m2}

\begin{figure}
\centering
\pandocbounded{\includegraphics[keepaspectratio,alt={base wave and harmonic m2}]{figs/waves_base_harmonic_m2.png}}
\caption{base wave and harmonic m2}
\end{figure}

\href{waves_base_harmonic_m2.svg}{SVG version}

\subsubsection{\texorpdfstring{Figure 8: Base wave and odd harmonic
\(m=3\)}{Figure 8: Base wave and odd harmonic m=3}}\label{figure-8-base-wave-and-odd-harmonic-m3}

\begin{figure}
\centering
\pandocbounded{\includegraphics[keepaspectratio,alt={base wave and harmonic m3}]{figs/waves_base_harmonic_m3.png}}
\caption{base wave and harmonic m3}
\end{figure}

\href{waves_base_harmonic_m3.svg}{SVG version}

The composite wave is

\[
\psi_{\mathrm{sum}}(\theta)
=
\cos\theta+\cos(m\theta)
\]

\subsubsection{\texorpdfstring{Figure 9: Composite wave
\(m=2\)}{Figure 9: Composite wave m=2}}\label{figure-9-composite-wave-m2}

\begin{figure}
\centering
\pandocbounded{\includegraphics[keepaspectratio,alt={composite waveform m2}]{figs/composite_wave_m2.png}}
\caption{composite waveform m2}
\end{figure}

\href{composite_wave_m2.svg}{SVG version}

\subsubsection{\texorpdfstring{Figure 10: Composite wave
\(m=3\)}{Figure 10: Composite wave m=3}}\label{figure-10-composite-wave-m3}

\begin{figure}
\centering
\pandocbounded{\includegraphics[keepaspectratio,alt={composite waveform m3}]{figs/composite_wave_m3.png}}
\caption{composite waveform m3}
\end{figure}

\href{composite_wave_m3.svg}{SVG version}

These waveform figures, too, are readout displays separate from the
state network. The construction of the complete network can be defined
without using the waveform figures.

\begin{center}\rule{0.5\linewidth}{0.5pt}\end{center}

\section{12. Complex-Vector Displays}\label{complex-vector-displays}

At each state \(k\), the two complex numbers

\[
z_1=U^k,
\]

\[
z_m=U^{mk}
\]

can be displayed as vectors on the unit circle.

\subsubsection{\texorpdfstring{Figure 11: \(m=2\): the two complex
vectors at all 4
states}{Figure 11: m=2: the two complex vectors at all 4 states}}\label{figure-11-m2-the-two-complex-vectors-at-all-4-states}

\begin{figure}
\centering
\pandocbounded{\includegraphics[keepaspectratio,alt={complex vectors all states m2}]{figs/complex_vectors_all_states_m2.png}}
\caption{complex vectors all states m2}
\end{figure}

\href{complex_vectors_all_states_m2.svg}{SVG version}

\subsubsection{\texorpdfstring{Figure 12: \(m=3\): the two complex
vectors at all 4
states}{Figure 12: m=3: the two complex vectors at all 4 states}}\label{figure-12-m3-the-two-complex-vectors-at-all-4-states}

\begin{figure}
\centering
\pandocbounded{\includegraphics[keepaspectratio,alt={complex vectors all states m3}]{figs/complex_vectors_all_states_m3.png}}
\caption{complex vectors all states m3}
\end{figure}

\href{complex_vectors_all_states_m3.svg}{SVG version}

For \(m=2\), the harmonic vector degenerates to the two directions
\(+1,-1\), whereas for \(m=3\) it cycles through four directions ---
this is directly visible.

Although not the discrete states themselves, figures confirming the
relative orientation of the two waves at the representative
continuous-phase position \(\theta=\pi/4\) are also kept.

\subsubsection{\texorpdfstring{Figure 13: Representative complex vectors
at
\(m=2,\ \theta=\pi/4\)}{Figure 13: Representative complex vectors at m=2,\textbackslash{} \textbackslash theta=\textbackslash pi/4}}\label{figure-13-representative-complex-vectors-at-m2-thetapi4}

\begin{figure}
\centering
\pandocbounded{\includegraphics[keepaspectratio,alt={complex vectors theta pi4 m2}]{figs/complex_vectors_theta_pi4_m2.png}}
\caption{complex vectors theta pi4 m2}
\end{figure}

\href{complex_vectors_theta_pi4_m2.svg}{SVG version}

\subsubsection{\texorpdfstring{Figure 14: Representative complex vectors
at
\(m=3,\ \theta=\pi/4\)}{Figure 14: Representative complex vectors at m=3,\textbackslash{} \textbackslash theta=\textbackslash pi/4}}\label{figure-14-representative-complex-vectors-at-m3-thetapi4}

\begin{figure}
\centering
\pandocbounded{\includegraphics[keepaspectratio,alt={complex vectors theta pi4 m3}]{figs/complex_vectors_theta_pi4_m3.png}}
\caption{complex vectors theta pi4 m3}
\end{figure}

\href{complex_vectors_theta_pi4_m3.svg}{SVG version}

Figures 13 and 14 are not the discrete \(k\) states themselves; they are
auxiliary figures on the continuous-phase display. This distinction is
maintained.

\begin{center}\rule{0.5\linewidth}{0.5pt}\end{center}

\section{13. Complete Expansion of the Squared
Sum}\label{complete-expansion-of-the-squared-sum}

\subsection{13.1 What is being squared and
summed}\label{what-is-being-squared-and-summed}

The squared sum in this section is not the squared sum of the 6 network
edge lengths \(d_{ij}\).

For the two wave components constituting each vertex state \(S_k\),

\[
z_1(k),
\qquad
z_m(k)
\]

first define, per state,

\[
Q_k
=
[z_1(k)]^2+[z_m(k)]^2
\]

and then compute, over all 4 states,

\[
Q_{\mathrm{loop}}
=
\sum_{k=0}^{3}Q_k
\]

This is therefore \textbf{the total squared sum of the two-component
complex states attached to the 4 vertices}.

Note that the conjugated norm sum at each state is

\[
|z_1(k)|^2+|z_m(k)|^2=1+1=2
\]

so over the 4 states

\[
\sum_{k=0}^{3}\left(|z_1(k)|^2+|z_m(k)|^2\right)=8
\]

which is constant regardless of \(m\) and hence carries no
discriminating power for this comparison. In contrast, the squared sum
without conjugation reduces to

\[
Q_{\mathrm{loop}}(m)
=
\sum_k U^{2k}+\sum_k U^{2mk}
\]

a \textbf{finite geometric series in doubled phase}. Its value can
therefore vary with the divisibility structure of the harmonic index
\(m\) and the cycle length \(K\). This is why this paper examines
\(Q_{\mathrm{loop}}\); it is not because probabilities or norms are
being identified with \(\sum z^2\).

\subsection{13.2 Expansion into real and imaginary
parts}\label{expansion-into-real-and-imaginary-parts}

Write a general complex number as

\[
z_j
=
a_j+i b_j
\]

Squaring,

\[
z_j^2
=
(a_j+i b_j)^2.
\]

Expanding the product,

\[
(a_j+i b_j)^2
=
a_j^2+2i a_j b_j+(i b_j)^2.
\]

From \(i^2=-1\),

\[
(i b_j)^2
=
i^2 b_j^2
=
-b_j^2.
\]

Therefore

\[
\boxed{
z_j^2
=
a_j^2-b_j^2+2i a_j b_j
}
\]

Summing over components,

\[
\sum_j z_j^2
=
\sum_j(a_j^2-b_j^2+2i a_jb_j).
\]

By linearity of the sum,

\[
\boxed{
\sum_j z_j^2
=
\sum_j a_j^2
-
\sum_j b_j^2
+
2i\sum_j a_jb_j
}
\]

Therefore, the real conditions for

\[
\sum_j z_j^2=0
\]

are

\[
\sum_j a_j^2
=
\sum_j b_j^2
\]

and

\[
\sum_j a_jb_j
=
0
\]

\subsection{\texorpdfstring{13.3 State-by-state computation for
\(m=2\)}{13.3 State-by-state computation for m=2}}\label{state-by-state-computation-for-m2}

\subsubsection{\texorpdfstring{\(k=0\)}{k=0}}\label{k0-2}

\[
S_0=(1,1).
\]

Therefore

\[
Q_0
=
1^2+1^2
=
2.
\]

\subsubsection{\texorpdfstring{\(k=1\)}{k=1}}\label{k1-2}

\[
S_1=(i,-1).
\]

From

\[
i^2=-1,
\]

\[
(-1)^2=1
\]

we get

\[
Q_1
=
-1+1
=
0.
\]

\subsubsection{\texorpdfstring{\(k=2\)}{k=2}}\label{k2-2}

\[
S_2=(-1,1).
\]

\[
Q_2
=
(-1)^2+1^2
=
1+1
=
2.
\]

\subsubsection{\texorpdfstring{\(k=3\)}{k=3}}\label{k3-2}

\[
S_3=(-i,-1).
\]

Since

\[
(-i)^2
=
(-1)^2i^2
=
-1,
\]

\[
(-1)^2=1
\]

we get

\[
Q_3
=
-1+1
=
0.
\]

Hence

\[
\boxed{
(Q_0,Q_1,Q_2,Q_3)
=
(2,0,2,0)
}
\]

and

\[
Q_{\mathrm{loop}}(2)
=
2+0+2+0
=
\boxed{4}.
\]

\subsection{\texorpdfstring{13.4 State-by-state computation for
\(m=3\)}{13.4 State-by-state computation for m=3}}\label{state-by-state-computation-for-m3}

\subsubsection{\texorpdfstring{\(k=0\)}{k=0}}\label{k0-3}

\[
S_0=(1,1)
\]

so

\[
Q_0=2.
\]

\subsubsection{\texorpdfstring{\(k=1\)}{k=1}}\label{k1-3}

\[
S_1=(i,-i).
\]

From

\[
i^2=-1,
\]

\[
(-i)^2=-1
\]

we get

\[
Q_1
=
-1-1
=
-2.
\]

\subsubsection{\texorpdfstring{\(k=2\)}{k=2}}\label{k2-3}

\[
S_2=(-1,-1)
\]

so

\[
Q_2
=
1+1
=
2.
\]

\subsubsection{\texorpdfstring{\(k=3\)}{k=3}}\label{k3-3}

\[
S_3=(-i,i).
\]

\[
Q_3
=
(-i)^2+i^2
=
-1-1
=
-2.
\]

Hence

\[
\boxed{
(Q_0,Q_1,Q_2,Q_3)
=
(2,-2,2,-2)
}
\]

and

\[
Q_{\mathrm{loop}}(3)
=
2-2+2-2
=
\boxed{0}.
\]

\subsection{13.5 Aggregation of the squared real and imaginary
components}\label{aggregation-of-the-squared-real-and-imaginary-components}

Collecting all 4 states and both components: for \(m=2\),

\[
\sum a^2=6,
\]

\[
\sum b^2=2,
\]

\[
\sum ab=0.
\]

Therefore

\[
\sum z^2
=
6-2+2i\times0
=
4.
\]

For \(m=3\),

\[
\sum a^2=4,
\]

\[
\sum b^2=4,
\]

\[
\sum ab=0.
\]

Therefore

\[
\sum z^2
=
4-4+2i\times0
=
0.
\]

On the other hand, the linear sums of the complex numbers themselves
vanish in both systems:

\[
\sum z_1(k)=0,
\]

for \(m=2\)

\[
\sum z_2(k)
=
1-1+1-1
=
0,
\]

for \(m=3\)

\[
\sum z_3(k)
=
1-i-1+i
=
0.
\]

Therefore ``zero linear sum'' and ``zero squared sum'' are distinct
conditions, and \(m=2\) satisfies the former but not the latter.

\begin{center}\rule{0.5\linewidth}{0.5pt}\end{center}

\section{\texorpdfstring{14. Exact Derivation of the Parity Selection
Rule at
\(K=4\)}{14. Exact Derivation of the Parity Selection Rule at K=4}}\label{exact-derivation-of-the-parity-selection-rule-at-k4}

To check whether the difference between \(m=2,3\) is an accidental
numerical instance, we derive the total squared sum at \(K=4\) for
arbitrary integer \(m\).

By definition,

\[
Q_{\mathrm{loop}}(m)
=
\sum_{k=0}^{3}
\left[(U^k)^2+(U^{mk})^2\right].
\]

From the law of exponents,

\[
(U^k)^2
=
U^{2k},
\]

\[
(U^{mk})^2
=
U^{2mk}.
\]

Therefore

\[
Q_{\mathrm{loop}}(m)
=
\sum_{k=0}^{3}U^{2k}
+
\sum_{k=0}^{3}U^{2mk}.
\]

At \(K=4\), \(U=i\), so

\[
U^2=i^2=-1.
\]

The first term is

\[
\sum_{k=0}^{3}U^{2k}
=
\sum_{k=0}^{3}(-1)^k
\]

which is

\[
=1-1+1-1
=0.
\]

For the second term,

\[
U^{2m}
=(i^2)^m
=(-1)^m.
\]

Hence

\[
\sum_{k=0}^{3}U^{2mk}
=
\sum_{k=0}^{3}\left[(-1)^m\right]^k.
\]

\subsubsection{\texorpdfstring{For even
\(m\)}{For even m}}\label{for-even-m}

\[
(-1)^m=1
\]

so

\[
\sum_{k=0}^{3}\left[(-1)^m\right]^k
=
1+1+1+1
=4.
\]

Therefore

\[
Q_{\mathrm{loop}}(m)
=0+4
=4.
\]

\subsubsection{\texorpdfstring{For odd
\(m\)}{For odd m}}\label{for-odd-m}

\[
(-1)^m=-1
\]

so

\[
\sum_{k=0}^{3}\left[(-1)^m\right]^k
=
1-1+1-1
=0.
\]

Therefore

\[
Q_{\mathrm{loop}}(m)
=0+0
=0.
\]

Hence, at \(K=4\), exactly

\[
\boxed{
Q_{\mathrm{loop}}(m)
=
\begin{cases}
4,&m\ \mathrm{even},\\
0,&m\ \mathrm{odd}
\end{cases}
}
\]

Therefore \(m=2\) and \(m=3\) are not merely two individual examples;
they are the minimal nontrivial even and odd representatives at \(K=4\).

This parity selection rule at \(K=4\) can be written as a corollary of
the following general theorem. For the primitive \(K\)-th root

\[
U=e^{2\pi i/K}
\]

define

\[
Q_{\mathrm{loop}}(m;K)
=
\sum_{k=0}^{K-1}
\left(U^{2k}+U^{2mk}\right)
\]

For the finite geometric series, for an integer \(r\),

\[
\sum_{k=0}^{K-1}U^{rk}
=
\begin{cases}
K,&K\mid r,\\
0,&K\nmid r
\end{cases}
\]

Indeed, if \(K\mid r\) then \(U^r=1\), so all \(K\) terms are 1. If
\(K\nmid r\) then \(U^r\ne1\) and \((U^r)^K=1\), hence

\[
\sum_{k=0}^{K-1}(U^r)^k
=
\frac{1-(U^r)^K}{1-U^r}
=
0
\]

Evaluating the two series separately, we obtain

\[
\boxed{
Q_{\mathrm{loop}}(m;K)
=
K\,\mathbf 1_{K\mid2}
+
K\,\mathbf 1_{K\mid2m}
}
\]

where \(\mathbf 1_{P}\) equals 1 when condition \(P\) holds and 0
otherwise.

At \(K=4\), since \(4\nmid2\), the first term vanishes and

\[
Q_{\mathrm{loop}}(m;4)
=
4\,\mathbf 1_{4\mid2m}.
\]

Furthermore,

\[
4\mid2m
\iff
2\mid m
\]

so the previously obtained

\[
Q_{\mathrm{loop}}(m;4)
=
\begin{cases}
4,&m\ \mathrm{even},\\
0,&m\ \mathrm{odd}
\end{cases}
\]

follows immediately. Therefore the parity difference at \(K=4\) is not
an isolated numerical instance but an exact corollary of divisibility
conditions on the finite cyclic group. For general \(K\), it is not
``parity'' itself but \(K\mid2\) and \(K\mid2m\) that govern the
classification. For example, at \(K=6,m=3\) we have \(6\mid2m\), so
\(Q_{\mathrm{loop}}=6\): an odd \(m\) that does not close to zero.

\begin{center}\rule{0.5\linewidth}{0.5pt}\end{center}

\subsection{\texorpdfstring{14.1 A Note on \(m=1\) versus \(m=3\):
Distinction via the Gram
Structure}{14.1 A Note on m=1 versus m=3: Distinction via the Gram Structure}}\label{a-note-on-m1-versus-m3-distinction-via-the-gram-structure}

At \(K=4\), if one looks only at the distance multiset, the minimal
harmonic orders, and \(Q_{\mathrm{loop}}\), then \(m=1\) and \(m=3\) are
not distinguished. This corresponds to \(U^3=\overline U\).

The additional relational invariant that provides this distinction is
the standard Hermitian inner product between states,

\[
\langle S_j,S_k\rangle
=
\overline{U^j}U^k
+
\overline{U^{mj}}U^{mk}
\]

For adjacent states,

\[
\langle S_k,S_{k+1}\rangle
=
\overline{U^k}U^{k+1}
+
\overline{U^{mk}}U^{m(k+1)}
=
U+U^m
\]

which is independent of \(k\). Substituting \(U=i\),

\[
m=1:\quad i+i=2i,
\]

\[
m=2:\quad i+i^2=-1+i,
\]

\[
m=3:\quad i+i^3=i-i=0.
\]

Therefore

\[
\boxed{
m=1,\ m=3\text{ are indistinguishable by the distance multiset but distinguishable by the complex Gram structure}
}
\]

In particular, for \(m=3\) all adjacent inner products vanish, and

\[
S_{k+2}
=
\left(U^2 U^k,\;U^6 U^{3k}\right)
=
-S_k
\]

holds. Hence the 4 states take the form

\[
\boxed{
S_0,S_1,S_2,S_3
=
\{v,w,-v,-w\},
\qquad
\langle v,w\rangle=0
}
\]

an orthogonal-pair \(\pm\) structure in \(\mathbb C^2\). This result
provides an invariant that refines the classification of complete
networks beyond the distance multiset. The comparison of \(m=2\) versus
\(m=3\) remains valid as the minimal parity comparison, and the results
of this subsection do not alter that conclusion.

\begin{center}\rule{0.5\linewidth}{0.5pt}\end{center}

\section{\texorpdfstring{15. \(U^K=I\) and Minimal
Orders}{15. U\^{}K=I and Minimal Orders}}\label{uki-and-minimal-orders}

\subsection{15.1 General order of the harmonic
generator}\label{general-order-of-the-harmonic-generator}

Let \(U\) have order \(K\), and let \(r\) be the order of \(U^m\).

By the definition of order,

\[
(U^m)^r=1.
\]

This is equivalent to

\[
U^{mr}=1
\]

Since the minimal order of \(U\) is \(K\),

\[
K\mid mr
\]

is required.

Setting \(g=\gcd(K,m)\) and decomposing

\[
K=gK',
\qquad
m=gm'
\]

with

\[
\gcd(K',m')=1,
\]

the condition \(K\mid mr\) becomes

\[
gK'\mid gm'r
\]

so, cancelling \(g\),

\[
K'\mid m'r.
\]

Since \(K'\) and \(m'\) are coprime, the minimal positive integer \(r\)
is

\[
r=K'.
\]

Therefore

\[
\boxed{
\operatorname{ord}(U^m)
=
\frac{K}{\gcd(K,m)}
}
\]

\subsection{\texorpdfstring{15.2
\(K=4,\ m=2\)}{15.2 K=4,\textbackslash{} m=2}}\label{k4-m2}

\[
\gcd(4,2)=2
\]

so

\[
\operatorname{ord}(U^2)
=
\frac{4}{2}
=2.
\]

Indeed,

\[
U^2=-1,
\]

\[
(U^2)^2=(-1)^2=1.
\]

\subsection{\texorpdfstring{15.3
\(K=4,\ m=3\)}{15.3 K=4,\textbackslash{} m=3}}\label{k4-m3}

\[
\gcd(4,3)=1
\]

so

\[
\operatorname{ord}(U^3)
=
\frac{4}{1}
=4.
\]

Indeed,

\[
U^3=-i,
\]

\[
(-i)^2=-1,
\]

\[
(-i)^3=i,
\]

\[
(-i)^4=1.
\]

\subsection{15.4 The cycle of the full two-wave
state}\label{the-cycle-of-the-full-two-wave-state}

Since the first component \(U^k\) itself maintains the primitive 4-state
cycle, even though the harmonic degenerates to 2 states at \(m=2\), the
full two-wave state pair \(S_k\) first returns after 4 steps.

Therefore,

\[
\boxed{
\operatorname{ord}(S_k)=4
\quad\text{for both }m=2,3
}
\]

The difference lies in the fact that

\[
\boxed{
\text{for }m=2\text{, one of the internal components degenerates to period 2 first}
}
\]

\begin{center}\rule{0.5\linewidth}{0.5pt}\end{center}

\section{16. Dynamic Displays: Only the Discrete State Moves on a Fixed
Network}\label{dynamic-displays-only-the-discrete-state-moves-on-a-fixed-network}

To prevent misunderstanding of the discreteness of this paper, dynamic
HTML and MP4 files were created.

The important point is that the network background and the waveforms
themselves are not redrawn at each step; only the marker showing the
current state moves discretely through

\[
k=0\to1\to2\to3\to0
\]

\subsection{16.1 Dynamic HTML}\label{dynamic-html}

\begin{itemize}
\tightlist
\item
  \href{K4_harmonic_discrete_phase_animation_20260915.html}{combined
  version: network + complex vectors + waveforms}
\item
  \href{K4_network_wave_discrete_m2.html}{\(m=2\): discrete-phase HTML
  with the original complete-network figure and base/harmonic figures as
  fixed background}
\item
  \href{K4_network_wave_discrete_m3.html}{\(m=3\): discrete-phase HTML
  with the original complete-network figure and base/harmonic figures as
  fixed background}
\item
  \href{K4_network_wave_index.html}{HTML index}
\end{itemize}

In the individual HTML files for \(m=2,3\), the original static figures
are not replaced; a red circle marking the current state and a red
vertical dotted line marking the phase position are overlaid on them.

\subsection{16.2 MP4}\label{mp4}

\begin{itemize}
\tightlist
\item
  \href{K4_harmonic_discrete_phase_m2_20260915.mp4}{\(m=2\) forward}
\item
  \href{K4_harmonic_discrete_phase_m3_20260915.mp4}{\(m=3\) forward}
\item
  \href{K4_harmonic_discrete_phase_m2_reverse_20260915.mp4}{\(m=2\)
  reverse}
\item
  \href{K4_harmonic_discrete_phase_m3_reverse_20260915.mp4}{\(m=3\)
  reverse}
\end{itemize}

The reverse versions serve to confirm visually that reading the cycle in
reverse order returns to the same finite state set. This is not a
physical claim of time-reversal symmetry.

\begin{center}\rule{0.5\linewidth}{0.5pt}\end{center}

\section{17. What This Paper Establishes
Exactly}\label{what-this-paper-establishes-exactly}

We summarize the results so far without admixture of hypotheses.

\subsection{17.1 Complete relational
structure}\label{complete-relational-structure}

For \(N=1, K=4\), the total number of vertices is

\[
P=4
\]

and the total number of relations is

\[
M_P=6.
\]

Since all 6 relations are kept, the network is the complete graph
\(K_4\).

\subsection{17.2 Finite cycle}\label{finite-cycle}

Using the primitive root

\[
U=i
\]

the states

\[
S_k=(U^k,U^{mk})
\]

satisfy

\[
S_{k+4}=S_k
\]

exactly.

\subsection{17.3 Internal order difference of even and odd
harmonics}\label{internal-order-difference-of-even-and-odd-harmonics}

\[
m=2:
\quad
\operatorname{ord}(U^2)=2,
\]

\[
m=3:
\quad
\operatorname{ord}(U^3)=4.
\]

\subsection{17.4 Complete-network geometric difference once a
state-space distance is
defined}\label{complete-network-geometric-difference-once-a-state-space-distance-is-defined}

Defining the \(\mathbb C^2\) Euclidean distance,

\[
m=2:
\quad
\{d_{ij}\}
=
\{2,2,\sqrt6,\sqrt6,\sqrt6,\sqrt6\},
\]

\[
m=3:
\quad
\{d_{ij}\}
=
\{2,2,2,2,2\sqrt2,2\sqrt2\}.
\]

Thus, with the same combinatorial \(K_4\) structure, the distance
geometry defined from the state labels differs.

\subsection{17.5 Difference in squared closure of the two-wave state
components}\label{difference-in-squared-closure-of-the-two-wave-state-components}

\[
m=2:
\quad
Q_{\mathrm{loop}}=4,
\]

\[
m=3:
\quad
Q_{\mathrm{loop}}=0.
\]

Furthermore, for \(K=4\) and arbitrary integer \(m\), we derived exactly
that

\[
Q_{\mathrm{loop}}(m)
=
4
\]

for even \(m\) and

\[
Q_{\mathrm{loop}}(m)
=
0
\]

for odd \(m\).

\begin{center}\rule{0.5\linewidth}{0.5pt}\end{center}

\section{18. What This Paper Does Not Yet
Establish}\label{what-this-paper-does-not-yet-establish}

The following do not follow automatically from the mathematical results
of this paper.

\begin{enumerate}
\def\labelenumi{\arabic{enumi}.}
\tightlist
\item
  That \(N=1\) is a real single photon.
\item
  That the base wave plus one harmonic is the complete internal degrees
  of freedom of a photon.
\item
  That \(k=0,1,2,3\) are four instants of physical time.
\item
  That the period of \(U^K=I\) is a physical time period of the
  universe.
\item
  That \(d_{ij}\) is any of real spatial distance, proper time, geodesic
  length, or causal distance.
\item
  That the MDS coordinates are physical space coordinates.
\item
  That the ``inflation-like rapid expansion'' of the preceding
  relational-wave system is identical to cosmological inflation.
\item
  That the squared-zero closure at \(m=3\) alone can select real
  particle states.
\end{enumerate}

Peer evaluation should therefore proceed in two stages.

The first stage is:

\[
\boxed{
\text{are the mathematical results correct as a finite complete-relation closed system?}
}
\]

The second stage is:

\[
\boxed{
\text{which physical quantities is it valid to map this closed system onto?}
}
\]

The validity of the first stage cannot be denied on the grounds that the
second stage is unsettled. Conversely, the mathematical validity of the
first stage alone cannot establish the physical identifications of the
second stage.

\begin{center}\rule{0.5\linewidth}{0.5pt}\end{center}

\section{19. Relation to Prior
Concepts}\label{relation-to-prior-concepts}

To the author's knowledge, no prior work adopts the same construction as
this paper --- completely relating the \(KN\) states as anonymous
vertices and closing them by \(U^K=I\). Concepts close to parts of the
problem, however, exist.

\subsection{19.1 Page--Wootters: descriptions that do not take external
time as
fundamental}\label{pagewootters-descriptions-that-do-not-take-external-time-as-fundamental}

Page and Wootters {[}E1{]} proposed treating the closed total system as
a stationary state and describing observed time evolution through
correlations with an internal clock.

The point of contact with this paper is that external time is not taken
as a fundamental variable. However, Page--Wootters splits the Hilbert
space into clock/system, and does not integrate all states into
anonymous vertices of the same kind forming a complete graph as this
paper does.

\subsection{19.2 Discrete history states}\label{discrete-history-states}

Boette and Rossignoli {[}E2{]} treat discrete system-time history states
for finite-dimensional clocks, embedding states at multiple instants
into one static structure.

The point of contact is holding finite discrete history states at once.
The difference is that this paper assumes no separate clock Hilbert
space; the state instants themselves become vertices and all relations
are kept.

\subsection{19.3 All-at-once / time-symmetric boundary
conditions}\label{all-at-once-time-symmetric-boundary-conditions}

Wharton {[}E3{]} argued against taking the generation of the future from
initial conditions as the only fundamental form, considering instead
constraints on entire histories from boundary conditions including
initial and final ones.

The point of contact is not privileging the ``sequentially generate the
future'' description. On the other hand, this paper does not start from
action principles or continuous-field boundary-value problems; it starts
from finite cyclic groups and complete relational networks.

\subsection{19.4 Causal sets}\label{causal-sets}

Bombelli, Lee, Meyer, and Sorkin {[}E4{]} proposed the causal set:
treating spacetime as a set of discrete events and causal relations
prior to any continuous manifold.

The point of contact is taking relational structure as fundamental
without presupposing continuous coordinates. However, the fundamental
relation of a causal set is a causal partial order, mathematically
different from the complete graph of this paper, which keeps every
vertex pair.

These are therefore cited not as derivational grounds but as
\textbf{conceptually neighboring studies}.

\begin{center}\rule{0.5\linewidth}{0.5pt}\end{center}

\section{20. The Meaning of the Connection to the Preceding Inflation
System}\label{the-meaning-of-the-connection-to-the-preceding-inflation-system}

This paper does not claim to reproduce the inflation-like
relational-wave dynamics of Figures 1--4 from the \(K_4\) system alone.

The connection is restricted to the following problem setting.

The preceding system used the update

\[
\text{state}(\tau)
\to
\text{state}(\tau+\Delta\tau)
\]

This paper is the minimal exactly closed system for examining whether
this can be replaced by a finite state set

\[
\{S_0,S_1,\ldots,S_{K-1}\}
\]

together with its complete relations

\[
K_{KN}
\]

The question we ultimately want to ask is therefore

\[
\boxed{
\text{can the orbit order be read from the complete relations alone, without an external update variable?}
}
\]

What was confirmed here is part of the necessary conditions for that:

\begin{itemize}
\tightlist
\item
  A finite complete-relation system with a cycle can be constructed.
\item
  With only internal states of base wave plus harmonic, the same 4
  vertices can carry different relational geometries depending on
  parity.
\item
  At \(K=4\), the harmonic orders and the squared closure have exact
  parity differences.
\end{itemize}

This alone does not derive physical time. But neither does ``turning the
time direction into vertices'' break the structure.

\begin{center}\rule{0.5\linewidth}{0.5pt}\end{center}

\section{21. Future Verification Tasks}\label{future-verification-tasks}

Tasks naturally following from this paper, separated into mathematical
and physical ones.

\subsection{21.1 Mathematical tasks}\label{mathematical-tasks}

\subsubsection{\texorpdfstring{(a) Extension of the complete
classification for general
\(K\)}{(a) Extension of the complete classification for general K}}\label{a-extension-of-the-complete-classification-for-general-k}

For the total squared sum, this paper classified exactly

\[
Q_{\mathrm{loop}}(m;K)
=
K\,\mathbf 1_{K\mid2}
+
K\,\mathbf 1_{K\mid2m}
\]

The next task is to combine this divisibility classification with the
minimal order of the harmonic component

\[
\operatorname{ord}(U^m)=\frac{K}{\gcd(K,m)}
\]

and the distance multisets and Gram invariants (the invariants provided
by the inter-state complex inner products \(\langle S_i,S_j\rangle\)) of
the complete network, to classify completely the equivalence classes of
general \((K,m)\).

\subsubsection{\texorpdfstring{(b) General
\(N>1,\ K>1\)}{(b) General N\textgreater1,\textbackslash{} K\textgreater1}}\label{b-general-n1-k1}

For the total number of vertices

\[
P=KN
\]

with all relations

\[
\frac{P(P-1)}2
\]

kept anonymously, investigate whether the original \(N\)-decomposition
and \(K\)-decomposition can be reconstructed from relational invariants.

This is important. If, given only \(KN\) anonymous vertices from
outside, the \(N\times K\) decomposition had to be supplied a priori,
the anonymity problem would persist in a different form.

\subsubsection{(c) Relational reconstruction of
order}\label{c-relational-reconstruction-of-order}

A complete graph has all relations, but that alone does not
automatically specify a ``next state''. In this paper, the order is
known through the phase cycle of \(S_k=(U^k,U^{mk})\).

In the future, it must be verified whether, after erasing the order
label \(k\), the cycle

\[
S_0\to S_1\to\cdots\to S_{K-1}\to S_0
\]

can be reconstructed from relational quantities alone.

This is the central falsifiable point of this research program.

At \(K=4\), the undirected cyclic order is reconstructible from the
distance structure, while reconstructing the direction is a separate
problem. Whether directional information can be obtained from additional
invariants including the complex inner products will be treated in a
separate paper.

\subsection{21.2 Physical tasks}\label{physical-tasks}

\subsubsection{(a) Photon identification}\label{a-photon-identification}

Whether the \(N=1\) system of base wave plus one harmonic can reproduce
photon polarization, helicity, energy--momentum, and quantum statistics
is unverified.

\subsubsection{(b) Length readout}\label{b-length-readout}

The \(d_{ij}\) of this paper is the standard norm on \(\mathbb C^2\). To
claim physical spacetime distance, the same distance law must be derived
independently from measurement procedures or invariants.

\subsubsection{(c) Returning to the inflation
system}\label{c-returning-to-the-inflation-system}

The final connection test is to reconstruct the measured trajectories of
the preceding \(N=40\) system as a \(KN\) complete-relation system
without external \(\tau\) updates, and check agreement with the
conventional trajectories.

\begin{center}\rule{0.5\linewidth}{0.5pt}\end{center}

\section{22. Reproducibility Materials}\label{reproducibility-materials}

All static figures, dynamic displays, raw data, and analysis programs of
this paper are stored in the same folder as this Markdown file or in the
reproducibility-verification folder directly below it.

\subsection{22.1 Static figure
generation}\label{static-figure-generation}

\begin{itemize}
\tightlist
\item
  \href{generate_all_K4_harmonic_figures_20260915.py}{generation program
  for all K4 figures}
\item
  \href{figs/K4_complete_network_m2.png}{\(m=2\) complete network PNG}
\item
  \href{K4_complete_network_m2.svg}{\(m=2\) complete network SVG}
\item
  \href{figs/K4_complete_network_m3.png}{\(m=3\) complete network PNG}
\item
  \href{K4_complete_network_m3.svg}{\(m=3\) complete network SVG}
\item
  \href{figs/waves_base_harmonic_m2.png}{\(m=2\) base + harmonic PNG}
\item
  \href{waves_base_harmonic_m2.svg}{\(m=2\) base + harmonic SVG}
\item
  \href{figs/waves_base_harmonic_m3.png}{\(m=3\) base + harmonic PNG}
\item
  \href{waves_base_harmonic_m3.svg}{\(m=3\) base + harmonic SVG}
\item
  \href{figs/composite_wave_m2.png}{\(m=2\) composite wave PNG}
\item
  \href{composite_wave_m2.svg}{\(m=2\) composite wave SVG}
\item
  \href{figs/composite_wave_m3.png}{\(m=3\) composite wave PNG}
\item
  \href{composite_wave_m3.svg}{\(m=3\) composite wave SVG}
\item
  \href{figs/complex_vectors_all_states_m2.png}{\(m=2\) all-states
  complex vectors PNG}
\item
  \href{complex_vectors_all_states_m2.svg}{\(m=2\) all-states complex
  vectors SVG}
\item
  \href{figs/complex_vectors_all_states_m3.png}{\(m=3\) all-states
  complex vectors PNG}
\item
  \href{complex_vectors_all_states_m3.svg}{\(m=3\) all-states complex
  vectors SVG}
\item
  \href{figs/complex_vectors_theta_pi4_m2.png}{\(m=2,\ \theta=\pi/4\)
  PNG}
\item
  \href{complex_vectors_theta_pi4_m2.svg}{\(m=2,\ \theta=\pi/4\) SVG}
\item
  \href{figs/complex_vectors_theta_pi4_m3.png}{\(m=3,\ \theta=\pi/4\)
  PNG}
\item
  \href{complex_vectors_theta_pi4_m3.svg}{\(m=3,\ \theta=\pi/4\) SVG}
\end{itemize}

\subsection{22.2 Complete-network edge
data}\label{complete-network-edge-data}

\begin{itemize}
\tightlist
\item
  \href{edges_m2.csv}{\(m=2\) 6-edge CSV}
\item
  \href{edges_m3.csv}{\(m=3\) 6-edge CSV}
\item
  \href{all_edges_m2_m3.csv}{\(m=2,3\) combined CSV}
\end{itemize}

\subsection{22.3 Raw data of the discrete complex
states}\label{raw-data-of-the-discrete-complex-states}

\begin{itemize}
\tightlist
\item
  \href{reproducibility_raw_readout_20260915/}{raw-data readout folder}
\item
  \href{reproducibility_raw_readout_20260915/raw_data/complex_states_animation_exact_i.csv}{exact-\(i\)
  discrete states}
\item
  \href{reproducibility_raw_readout_20260915/raw_data/complex_states_static_exp.csv}{static
  \(\exp(i\theta)\) states}
\item
  \href{reproducibility_raw_readout_20260915/raw_data/complex_states.csv}{original
  complex-state CSV}
\item
  \href{reproducibility_raw_readout_20260915/raw_data/RAW_DATA_MANIFEST.csv}{raw
  data manifest}
\item
  \href{reproducibility_raw_readout_20260915/raw_data/representative_vectors_theta_pi4.csv}{\(\theta=\pi/4\)
  representative vectors}
\end{itemize}

\subsection{22.4 Waveform sample data}\label{waveform-sample-data}

\begin{itemize}
\tightlist
\item
  \href{reproducibility_raw_readout_20260915/raw_data/wave_samples_static_2000_m2.csv}{\(m=2\)
  static 2000 samples}
\item
  \href{reproducibility_raw_readout_20260915/raw_data/wave_samples_static_2000_m3.csv}{\(m=3\)
  static 2000 samples}
\item
  \href{reproducibility_raw_readout_20260915/raw_data/wave_samples_html_800_m2.csv}{\(m=2\)
  HTML 800 samples}
\item
  \href{reproducibility_raw_readout_20260915/raw_data/wave_samples_html_800_m3.csv}{\(m=3\)
  HTML 800 samples}
\item
  \href{reproducibility_raw_readout_20260915/raw_data/wave_samples_mp4_1200_m2.csv}{\(m=2\)
  MP4 1200 samples}
\item
  \href{reproducibility_raw_readout_20260915/raw_data/wave_samples_mp4_1200_m3.csv}{\(m=3\)
  MP4 1200 samples}
\end{itemize}

\subsection{22.5 Self-consistency and squared-sum
analysis}\label{self-consistency-and-squared-sum-analysis}

\begin{itemize}
\tightlist
\item
  \href{reproducibility_raw_readout_20260915/self_consistency_square_sum_U4_20260915/analyze_K4_self_consistency_square_sum_20260915.py}{analysis
  program}
\item
  \href{reproducibility_raw_readout_20260915/self_consistency_square_sum_U4_20260915/K4_self_consistency_square_sum_by_state_20260915.csv}{per-state
  results}
\item
  \href{reproducibility_raw_readout_20260915/self_consistency_square_sum_U4_20260915/K4_self_consistency_square_sum_summary_20260915.csv}{summary
  results}
\item
  \href{reproducibility_raw_readout_20260915/self_consistency_square_sum_U4_20260915/K4_self_consistency_square_sum_analysis_20260915.md}{analysis
  report}
\end{itemize}

\subsection{22.6 Dynamic HTML / MP4}\label{dynamic-html-mp4}

\begin{itemize}
\tightlist
\item
  \href{K4_harmonic_discrete_phase_animation_20260915.html}{combined
  HTML}
\item
  \href{K4_network_wave_discrete_m2.html}{\(m=2\) fixed-background HTML}
\item
  \href{K4_network_wave_discrete_m3.html}{\(m=3\) fixed-background HTML}
\item
  \href{K4_network_wave_index.html}{HTML index}
\item
  \href{K4_harmonic_discrete_phase_m2_20260915.mp4}{\(m=2\) forward MP4}
\item
  \href{K4_harmonic_discrete_phase_m3_20260915.mp4}{\(m=3\) forward MP4}
\item
  \href{K4_harmonic_discrete_phase_m2_reverse_20260915.mp4}{\(m=2\)
  reverse MP4}
\item
  \href{K4_harmonic_discrete_phase_m3_reverse_20260915.mp4}{\(m=3\)
  reverse MP4}
\item
  \href{generate_K4_harmonic_discrete_phase_mp4_20260915.py}{MP4
  generation program}
\item
  \href{generate_K4_harmonic_discrete_phase_forward_reverse_mp4_20260915.py}{forward/reverse
  MP4 generation program}
\item
  \href{generate_k4_phase_html.py}{HTML generation program}
\item
  \href{generate_K4_harmonic_complex_raw_states_20260915.py}{raw
  complex-state generation program}
\end{itemize}

\subsection{22.7 Figures reproduced from the preceding inflation
system}\label{figures-reproduced-from-the-preceding-inflation-system}

\begin{itemize}
\tightlist
\item
  \href{figs/inflation_N40_phase_step0_cross_en_20260915.png}{N=40 step0
  cross-shaped bands (English)}
\item
  \href{figs/inflation_N40_phase_step500_ring_en_20260915.png}{N=40
  step500 ring (English)}
\item
  \href{figs/inflation_mexican_hat_SSB_reconstruction_20260915.png}{Mexican-hat
  / SSB reconstruction}
\item
  \href{figs/inflation_N40_relation_wave_3D_step500_en_20260915.png}{N=40
  relation-wave 3D evolution (English)}
\end{itemize}

These four figures were placed in the same folder so that the background
of the preceding inflation system {[}K1{]} can be understood within this
paper. The English versions of Figures 1, 2, and 4 were faithfully
reconstructed from the canonical data
(\href{english_figures_inflation_en_20260915/}{reconstruction package}).

\begin{center}\rule{0.5\linewidth}{0.5pt}\end{center}

\section{23. Conclusion}\label{conclusion}

This paper proposed a construction that does not treat the time
direction as a special external update parameter: distinct phase states
themselves are treated on an equal footing with ordinary state vertices,
and all states are completely related. The minimal exactly closed system
\(N=1, K=4\) was analyzed completely.

The total number of vertices is

\[
P=KN
\]

and the total number of relations is

\[
M_P
=
\frac{KN(KN-1)}2
\]

At \(N=1, K=4\), this yields the complete graph \(K_4\) with 4 vertices
and 6 relations. Using the primitive root

\[
U=e^{2\pi i/4}=i
\]

and defining

\[
S_k=(U^k,U^{mk})
\]

an exact 4-state cycle was obtained without using external continuous
time.

Comparing the minimal nontrivial even harmonic \(m=2\) and odd harmonic
\(m=3\),

\[
m=2:
\quad
\operatorname{ord}(U^2)=2,
\]

\[
m=3:
\quad
\operatorname{ord}(U^3)=4,
\]

is the internal order difference.

Furthermore, when the \(\mathbb C^2\) Euclidean distance is defined as
an analytical readout, the distance multisets of the 6 edges differ:

\[
m=2:
\quad
\{2,2,\sqrt6,\sqrt6,\sqrt6,\sqrt6\},
\]

\[
m=3:
\quad
\{2,2,2,2,2\sqrt2,2\sqrt2\}
\]

Moreover, the total squared sum of the two-wave state components over
all 4 states is

\[
Q_{\mathrm{loop}}(2)=4,
\]

\[
Q_{\mathrm{loop}}(3)=0.
\]

And for \(K=4\) and arbitrary integer \(m\), we obtained the exact
parity selection rule

\[
\boxed{
Q_{\mathrm{loop}}(m)
=
\begin{cases}
4,&m\ \mathrm{even},\\
0,&m\ \mathrm{odd}
\end{cases}
}
\]

These claims require no photon identification, no physical-time
identification, and no physical-length readout. Conversely, reading
\(N=1\) as one photon, reading \(k\) as past/present/future, and reading
\(d_{ij}\) as physical spacetime distance were explicitly reserved as
hypotheses.

The central result of this paper is therefore:

\[
\boxed{
\text{without assuming external time evolution, a minimal exactly closed complete-relation system of finite phase states as vertices can be constructed,}
\\
\text{and the difference in its internal harmonic structure produces exact differences in orders, state-space distance geometry, and squared closure}
}
\]

shown while omitting no state and no relation.

The decisive next task is to verify whether the cycle order itself can
be reconstructed from the complete relational data alone, with the order
label \(k\) removed. If this is possible, a concrete mathematical bridge
is obtained from the description ``evolving states in time'' to the
description ``reading time order from all state relations''.

\begin{center}\rule{0.5\linewidth}{0.5pt}\end{center}

\section{References}\label{references}

\subsection{Self-citation}\label{self-citation}

\textbf{{[}K1{]}} Noriaki Kihara, ``The Mechanism of Inflation-like
Rapid Expansion in Self-Consistent Closed Systems of Relational Waves'',
2026. Version DOI:
\href{https://doi.org/10.5281/zenodo.22176949}{10.5281/zenodo.22176949};
Concept DOI:
\href{https://doi.org/10.5281/zenodo.22112008}{10.5281/zenodo.22112008}.

Used only as the source of the \(N\)-body complete relational-wave
system, the 90-degree high-symmetry floor, the rapid transverse growth,
and the quasi-stable ring that motivate this paper. The \(KN\)
vertexization, the \(K_4\) analysis, and the parity squared-closure
derivation of this paper are completed independently within this paper.

\subsection{External references: conceptually neighboring
studies}\label{external-references-conceptually-neighboring-studies}

\textbf{{[}E1{]}} D. N. Page and W. K. Wootters, ``Evolution without
evolution: Dynamics described by stationary observables,''
\emph{Physical Review D} \textbf{27}, 2885 (1983). DOI:
\href{https://doi.org/10.1103/PhysRevD.27.2885}{10.1103/PhysRevD.27.2885}.

\textbf{{[}E2{]}} A. Boette and R. Rossignoli, ``History states of
systems and operators,'' \emph{Physical Review A} \textbf{98}, 032108
(2018). DOI:
\href{https://doi.org/10.1103/PhysRevA.98.032108}{10.1103/PhysRevA.98.032108}.

\textbf{{[}E3{]}} K. Wharton, ``Time-Symmetric Boundary Conditions and
Quantum Foundations,'' \emph{Symmetry} \textbf{2}, 272--283 (2010). DOI:
\href{https://doi.org/10.3390/sym2010272}{10.3390/sym2010272}.

\textbf{{[}E4{]}} L. Bombelli, J. Lee, D. Meyer, and R. D. Sorkin,
``Space-time as a causal set,'' \emph{Physical Review Letters}
\textbf{59}, 521--524 (1987). DOI:
\href{https://doi.org/10.1103/PhysRevLett.59.521}{10.1103/PhysRevLett.59.521}.

\begin{center}\rule{0.5\linewidth}{0.5pt}\end{center}

\section{Appendix A: Minimal Table of
Results}\label{appendix-a-minimal-table-of-results}

{\def\LTcaptype{none} % do not increment counter
\begin{longtable}[]{@{}
  >{\raggedright\arraybackslash}p{(\linewidth - 4\tabcolsep) * \real{0.2727}}
  >{\raggedleft\arraybackslash}p{(\linewidth - 4\tabcolsep) * \real{0.3636}}
  >{\raggedleft\arraybackslash}p{(\linewidth - 4\tabcolsep) * \real{0.3636}}@{}}
\toprule\noalign{}
\begin{minipage}[b]{\linewidth}\raggedright
Item
\end{minipage} & \begin{minipage}[b]{\linewidth}\raggedleft
\(m=2\) even harmonic
\end{minipage} & \begin{minipage}[b]{\linewidth}\raggedleft
\(m=3\) odd harmonic
\end{minipage} \\
\midrule\noalign{}
\endhead
\bottomrule\noalign{}
\endlastfoot
Base generator & \(i\) & \(i\) \\
Minimal order of base & 4 & 4 \\
Harmonic generator & \(-1\) & \(-i\) \\
Minimal order of harmonic & 2 & 4 \\
Minimal loop of full two-wave state & 4 & 4 \\
Per-state \(Q_k\) & \([2,0,2,0]\) & \([2,-2,2,-2]\) \\
\(\sum a^2\) & 6 & 4 \\
\(\sum b^2\) & 2 & 4 \\
\(\sum ab\) & 0 & 0 \\
\(Q_{\mathrm{loop}}=\sum z^2\) & 4 & 0 \\
Zero squared closure & NO & YES \\
Short edges of the \(\mathbb C^2\) distance & 2 edges × 2 & 4 edges ×
2 \\
Long edges of the \(\mathbb C^2\) distance & 4 edges × \(\sqrt6\) & 2
edges × \(2\sqrt2\) \\
\end{longtable}
}

\begin{center}\rule{0.5\linewidth}{0.5pt}\end{center}

\section{Appendix B: Strict Usage of
Terms}\label{appendix-b-strict-usage-of-terms}

\textbf{Complete network:} the combinatorial structure keeping every
vertex pair as a relation. No physical distance is assumed.\\
\textbf{State-space distance:} the \(d_{ij}\) defined in this paper as
the standard norm on \(\mathbb C^2\). Not physical distance.\\
\textbf{MDS coordinates:} embedding coordinates for visualizing the
distance matrix. Not physical coordinates.\\
\textbf{Discrete phase states:} \(k=0,1,2,3\). They do not mean
continuous time.\\
\textbf{Waveform figures:} readout figures displaying \(\cos\theta\) and
\(\cos(m\theta)\) continuously. The actual network states are 4 discrete
points.\\
\textbf{Algebraic self-consistency:} the internal consistency of the
state generation rule, the harmonic relation, and the cycle as defined
in this paper. Not synonymous with the fixed-point condition of the
preceding nonlinear dynamics.\\
\textbf{Zero squared closure:} in this paper, \(Q_{\mathrm{loop}}=0\)
for the two-wave state components at the 4 vertices. Not the squared sum
of the 6 edge lengths.\\
\textbf{Inflation-like:} refers to the rapid order-parameter growth
observed in the preceding relational-wave system. No identity with
cosmological inflation is claimed.\\
\textbf{Photon model:} a candidate physical interpretation of the
\(N=1\) closed system. Not a precondition of the mathematical network
analysis.

\end{document}
